% !TeX program = tectonic
\documentclass[11pt,a4paper]{article}
\usepackage{fontspec}
\usepackage[russian]{babel}
\babelfont{rm}[Language=Default]{Liberation Serif}
\babelfont[Language=Default]{sf}{Carlito}
\babelfont[Language=Default]{tt}{DejaVu Sans Mono}
\usepackage{amsmath,amssymb,amsthm}
\usepackage{geometry}
\geometry{a4paper, top=2.4cm, bottom=2.4cm, left=2.4cm, right=2.4cm}
\usepackage{booktabs,tabularx,enumitem,xcolor,tcolorbox}
\tcbuselibrary{breakable,skins}
\definecolor{linkblue}{HTML}{1F4E79}
\definecolor{statgreen}{HTML}{2F6B3A}
\definecolor{goldacc}{HTML}{8B7E5A}
\usepackage[colorlinks=true,linkcolor=linkblue,urlcolor=linkblue,unicode=true]{hyperref}
\theoremstyle{plain}
\newtheorem{theorem}{Теорема}
\newtheorem{lemma}[theorem]{Лемма}
\newtheorem{proposition}[theorem]{Предложение}
\newtheorem{corollary}[theorem]{Следствие}
\theoremstyle{definition}
\newtheorem{definition}[theorem]{Определение}
\theoremstyle{remark}
\newtheorem{remark}[theorem]{Замечание}
\newtcolorbox{statusbox}[1][]{colback=statgreen!5,colframe=statgreen!75,
  title={\textbf{Сводка доказанного:} #1},fonttitle=\small,boxrule=0.7pt,arc=2pt,breakable}
\newtcolorbox{databox}[1][]{colback=goldacc!6,colframe=goldacc!80,
  title={\textbf{Данные протокола:} #1},fonttitle=\small,boxrule=0.7pt,arc=2pt,breakable}
\newtcolorbox{verifybox}[1][]{colback=linkblue!5,colframe=linkblue!80,
  title={\textbf{Верификация:} #1},fonttitle=\small,boxrule=0.7pt,arc=2pt,breakable}
\widowpenalty=10000
\clubpenalty=10000
\setlength{\parskip}{0.35em}
\newcommand{\CC}{\mathbb{C}}
\newcommand{\RR}{\mathbb{R}}
\newcommand{\ZZ}{\mathbb{Z}}
\newcommand{\QQ}{\mathbb{Q}}
\newcommand{\Ga}{\Gamma}
\newcommand{\Bfun}{\mathrm{B}}
\newcommand{\im}{\operatorname{Im}}
\newcommand{\re}{\operatorname{Re}}
\newcommand{\lcm}{\operatorname{lcm}}
\newcommand{\SNF}{\operatorname{SNF}}
\newcommand{\Jac}{\operatorname{Jac}}
\newcommand{\rank}{\operatorname{rank}}
\newcommand{\Ind}{\operatorname{Index}}
\newcommand{\disc}{\operatorname{disc}}
\begin{document}
\begin{center}
{\LARGE\bfseries Сертификат C: μ₄-эквивариантность поправок}\\[0.4em]
{\large 22 = 1 + 7 + 7 + 7 и регулярное представление}\\[0.6em]
{\normalsize Исаев Исхак Хамзатович}\\[0.2em]
{\small Программа <<Динамический принцип>> --- лаборатория \texttt{hodge-laboratory}}\\[0.15em]
{\small Отдельная монография теоремы · Монография 11/16}
\end{center}
\vspace{0.4em}
{\color{linkblue}\hrule height 0.8pt}
\vspace{0.8em}

\section{Постановка}

Сертификат C устанавливает естественность поправок: тройка
$(7,22,\lambda_1)$ согласована с представлением $\mu_4$ ---
группы четверть-оборотов, присутствующей на всех трёх стендах.

\begin{theorem}[C1 --- изотипика K3]\label{th:C1stand}
В $H^2$ Фермат-квартики изотипное разложение по характерам $\mu_4$
имеет вид $22=1+7+7+7$. Три независимых пути подсчёта (прямые
орбиты, спектр перестановочного представления, точные кратности
характеров) дают одинаковый ответ.
\end{theorem}

\begin{proof}
Действие $(\mu_4)^2$ на $48$ прямых имеет $12$ орбит.
Перестановочное представление раскладывается по характерам;
проекционные операторы $\Pi_\chi=\frac14\sum_k i^{-k}R^k$
выделяют изотипные части. Прямой подсчёт орбит даёт $(1,7,7,7)$;
спектр перестановочного представления даёт те же кратности как
размерности собственных пространств; спаривание с точными
характерами прямых --- третий путь. Совпадение трёх путей исключает
вычислительные совпадения.
\end{proof}

\begin{theorem}[C2 --- тор, регулярное представление]\label{th:C2stand}
На торе оператор поправки $\lambda_1$ коммутирует с $\mu_4$,
характеристический многочлен равен $x^4-1$ --- регулярному
представлению, кратность $\lambda_1$ равна $4$, континуум значений
$\lambda_1=4\pi^2$, тройка $(4,1,4\pi^2)$ согласована.
\end{theorem}

\begin{proof}
Коммутация --- эквивариантность лапласиана относительно поворотов;
собственные пространства инвариантны и раскладываются по характерам
$\mu_4$. Многочлен $x^4-1$ имеет простые корни $\{1,-1,i,-i\}$ ---
каждый характер ровно один раз: регулярное представление. Значение
$4\pi^2$ реализуется на каждом характере --- континуум на решётке
моментов.
\end{proof}

\begin{databox}[гептада Клейна]
$\Phi_7$ --- циклотомический многочлен; $\Delta_{49a1}=-343=-7^3$;
$j=-3375=-15^3$; CM-порядок $\ZZ[(1+\sqrt{-7})/2]$; $h(-7)=1$;
$\tau$ сертифицирован. Гептада $\sqrt{-7}$ --- точное целочисленное
окружение фазы $\pi/7$.
$\mu_4$-стенд: аномалия $=i$; $42$ элемента порядка $4$ в
перестановочном представлении размерности $32$; $\mu_4$ на всех
трёх стендах. Сетка: нарушений симметричной кодировки $0$;
нарушений $(7,22)$: $1023/1024$.
\end{databox}

\begin{statusbox}[итог]
Сертификат C принят ($21/21$ проверка): изотипика $22=1+7+7+7$
подтверждена тремя путями; тор даёт регулярное представление;
семёрка получает точное арифметическое лицо.
\end{statusbox}

\begin{verifybox}[как воспроизвести]
\texttt{python3 laboratory.py} $\to$ <<Сертификаты $\to$ C>>:
проекционные операторы, орбиты, спектр --- всё точно.
\end{verifybox}

\vfill
{\color{linkblue}\hrule height 0.6pt}
\vspace{0.5em}
{\small\color{gray}
Лицензия: индивидуальная исключительная лицензия Исаева Исхака Хамзатовича.
Все права на вычисления, формулы и тексты принадлежат автору.
Верификация: \texttt{python3 laboratory.py} (пункт <<Стенды>>/<<Сертификаты>>),
Lean: \texttt{verification/lean/HodgeLaboratory.lean}.
}
\end{document}
